\documentclass[11pt]{article}
\usepackage[T1]{fontenc}
\usepackage{lmodern}
\usepackage[margin=1in]{geometry}
\usepackage{amsmath,amssymb,amsthm,mathtools}
\usepackage{microtype}
\usepackage[numbers,sort&compress]{natbib}
\usepackage{booktabs}
\usepackage[hidelinks]{hyperref}
\hypersetup{pdftitle={Polynomial-time counting of two-level transpose fibres},
  pdfauthor={Xiaojian Zeng},pdfsubject={Exact enumeration of two-level transpose fibres}}
\newtheorem{theorem}{Theorem}
\newtheorem{proposition}[theorem]{Proposition}
\newtheorem{lemma}[theorem]{Lemma}
\theoremstyle{remark}
\newtheorem{remark}[theorem]{Remark}
\newcommand{\Z}{\mathbb Z}
\newcommand{\R}{\mathbb R}


\newcommand{\card}[1]{\lvert #1\rvert}
\DeclareMathOperator{\degp}{deg}
\title{Polynomial-time counting of two-level transpose fibres}
\author{Xiaojian Zeng\\[3pt]
\small PBC School of Finance, Tsinghua University, 100190\\
\small\texttt{xiaojianzeng@foxmail.com}}
\date{}
\begin{document}
\maketitle
\begin{abstract}
Barrese, Loehr, Remmel and Sagan asked how many Ferrers boards have a
specified singleton image under the $m$-level transpose. We give an exact
algorithm for the entire $m=2$ case. A row-pair bijection identifies each
fibre with nonnegative integer vectors satisfying adjacent-sum bounds and
one terminal bound. For $t$ coordinates and arbitrary nonnegative
capacities of maximum $A$, a discrete polynomial algorithm uses $O(t^4)$
integer operations on $O(t\log(A+t+2))$-bit integers and stores $O(t^2)$
integers. Thus $O(t^6\log^2(A+t+2))$ elementary bit operations suffice.
The key is a representation with at most $i$ integer-interval polynomial
pieces at layer $i$, normalized by samples inside each interval.
The original arbitrary-$m$ question remains unresolved here.
\end{abstract}

\section{The fibre problem and the result}

A Ferrers shape is a finite weakly increasing sequence of positive
column heights $B=(b_1,\ldots,b_n)$. The empty sequence is allowed.
We identify representations that differ only by leading zero columns,
as in Barrese's convention \cite[Section~3]{Barrese2021}.
For an integer $m\geq1$, level $p$ consists of rows
$(p-1)m+1,\ldots,pm$, counted from the bottom. Its cell total is
\[
 L_p(B)=\sum_{j=1}^n\min\{m,\max\{0,b_j-(p-1)m\}\}.
\]
If $B$ has $t$ nonempty levels, its \emph{$m$-level transpose} is
\[
 \ell_m(B)=(L_t(B),\ldots,L_1(B)),
 \qquad \ell_m(\varnothing)=\varnothing.
\]
A shape is \emph{$m$-singleton} if at most one column has height strictly
between $(p-1)m$ and $pm$, for each $p$. The level transpose maps Ferrers
shapes onto singleton shapes and is an involution on the singleton
shapes \cite[Definition~4 and the following paragraph]{BLRS2016}.
For a singleton target $C$, write
\[
 F_m(C)=\card{\{B:\ell_m(B)=C\}}.
\]
The width of $B$ is unrestricted. This set is finite because its members
have area $\sum_j c_j$.

Barrese, Loehr, Remmel and Sagan posed this counting question in
\cite[Section~6.2, p.~33]{BLRS2016}. Their earlier formula
\cite[Theorem~5.2]{BLRS2014} counts singleton boards in a whole
$m$-level rook-equivalence class. A fibre fixes one transpose image and
includes nonsingleton boards, so it is a different enumeration problem.
We evaluate every fibre for $m=2$. The arbitrary-$m$ question and the
broader enumeration of whole equivalence classes remain unresolved here.

For a nonnegative integer vector $a=(a_1,\ldots,a_t)$, define
\begin{equation}\label{eq:polytope}
 \mathcal D(a)=\{d\in\Z_{\geq0}^t:
 d_i+d_{i+1}\leq a_i\ (1\leq i<t),\quad d_t\leq a_t\}.
\end{equation}
The empty vector has one solution. All logarithms below are to base two.

\begin{theorem}\label{thm:main}
For every $t\geq1$ and $a\in\Z_{\geq0}^t$, set $A=\max_i a_i$.
The number $\card{\mathcal D(a)}$ can be computed exactly with
$O(t^4)$ integer arithmetic operations on integers of
$O(t\log(A+t+2))$ bits, using $O(t^2)$ stored integers.
In particular, $O(t^6\log^2(A+t+2))$ elementary bit operations suffice.
At layer $i$, the algorithm represents its count function on at most
$i$ integer intervals, each by a polynomial of degree at most $i-1$.

If $C=(c_1,\ldots,c_t)$ is a nonempty $2$-singleton shape, put
\begin{equation}\label{eq:capacity}
 s_i=c_{t+1-i},\qquad s_{t+1}=0,\qquad
 a_i=\left\lfloor \frac{s_i}{2}\right\rfloor-\left\lceil \frac{s_{i+1}}{2}\right\rceil.
\end{equation}
Then $a_i\geq0$ and $F_2(C)=\card{\mathcal D(a)}$ by an explicit
bijection. Consequently $F_2(C)$ is computable in time polynomial in
the binary input length
\[
 L=\log(t+1)+\sum_{i=1}^t\log(c_i+1).
\]
The empty target has $F_2(\varnothing)=1$.
\end{theorem}

The quantitative issue is dependence on the \emph{bit lengths} of the
capacities. A table indexed by every possible coordinate value can be
exponentially large in those lengths. The theorem replaces such a table
by a bounded number of univariate pieces; it does not require a fixed
dimension or a repeated pattern in $a$.

\paragraph{Relation to earlier counting methods.}
For a common adjacent-sum capacity $q$, B\'ona, Ju and Yoshida
\cite[Section~2.1]{BJY2007} express weighted-path counts through a
$(q+1)$-square transfer matrix. Jiang, Yang and Zhong
\cite[Theorem~1.1 and equation~(4)]{JYZ2026} extend this common-capacity
model to blocks, compressing the transfer matrix to $q+1$ height states.
Jiang, Wen and Zhong
\cite[Section~2.1]{JWZ2026} use transfer matrices for bounds alternating
between $q$ and $q+1$. Cigler and Krattenthaler
\cite[Section~3, Theorem~4]{CK2024} enumerate alternating sequences in
a common interval with specified endpoints, and develop weighted and
tableau extensions. These works provide formulas for structured
capacities; the present bound holds for every input vector $a$.

There is also a bounded-poset interpretation. With
$\beta_1=0$, $\beta_{i+1}=a_i-\beta_i$ and
$y_i=(-1)^{i-1}(d_i-\beta_i)$, the edge inequalities become
$y_1\leq y_2\geq y_3\leq\cdots$. The individual bounds on $d_i$
become integer lower and upper bounds on $y_i$. Attaching marked lower
and upper elements realizes a marked order polytope.
Jochemko and Sanyal \cite[Theorem~2.6]{JS2014} prove piecewise
polynomiality in its markings. Assaf and Bergeron
\cite[Corollary~3.15]{AB2020} give a linear-extension decomposition
for bounded $\mathcal P$-partitions. Neither statement by itself bounds
the computation polynomially in the variable poset size. Our algorithm
works at one input and maintains at most $i$ univariate pieces at layer $i$.

Symbolic inference already uses piecewise polynomials and interpolation:
Zhe Zeng and Van den Broeck \cite[Sections~4.1--4.3]{ZVB2020} develop
these methods for integration over real variables. Here the operation is
discrete prefix summation, with a separate coefficient-size analysis.
The distinction between arithmetic operations and bit complexity is also
explicit in Zijian Zeng's treatment of polynomial evaluation and
interpolation \cite{Zeng2026}. Our proof supplies both
bounds for consecutive integer sampling nodes.

The contribution is the combination of the fibre bijection and the
uniform interval and bit bounds. In the primary statements inspected
through 13 September 2026, we found no equivalent binary-input fibre
theorem. This is a bounded literature comparison; it does not exclude
all possible specializations of earlier symbolic or tableau methods.

\section{The row-pair bijection}

For $m=2$, the singleton condition says exactly that odd column heights
do not repeat. If $u\leq v$ are consecutive target heights, then
$\left\lceil u/2\right\rceil\leq\left\lfloor v/2\right\rfloor$: failure would require $u=v$ odd.
Thus the capacities in \eqref{eq:capacity} are nonnegative.

\begin{proposition}\label{prop:bijection}
With the notation of \eqref{eq:capacity}, the fibre of $C$ is in bijection
with $\mathcal D(a)$. For $d\in\mathcal D(a)$, its preimage has row lengths
\begin{equation}\label{eq:rows}
 r_{2i-1}=\left\lceil \frac{s_i}{2}\right\rceil+d_i,
 \qquad r_{2i}=\left\lfloor \frac{s_i}{2}\right\rfloor-d_i
 \quad(1\leq i\leq t).
\end{equation}
\end{proposition}
\begin{proof}
Given a preimage $B$, let $r_j$ be its number of columns of height at
least $j$, padding with zero rows up to $2t$. These row lengths are
weakly decreasing and nonnegative, and
$r_{2i-1}+r_{2i}=s_i$. Setting
$d_i=\left\lfloor s_i/2\right\rfloor-r_{2i}$ gives \eqref{eq:rows}.
The within-pair order is equivalent to $d_i\geq0$; the between-pair
order $r_{2i}\geq r_{2i+1}$ is equivalent to
$d_i+d_{i+1}\leq a_i$. The last row is nonnegative precisely when
$d_t\leq a_t$.

Conversely these inequalities make every row in \eqref{eq:rows}
weakly decreasing, with the last row nonnegative. Since $s_t>0$,
there are exactly $t$ nonempty levels. Conjugating the rows gives a
unique shape with positive column heights and the prescribed level
totals. These constructions are inverse. Its width is $r_1$ and is
allowed to vary. Finally only the empty shape has empty transpose.
\end{proof}

For example, $C=(2,4)$ gives $a=(1,1)$. The three vectors
$(0,0),(1,0),(0,1)$ yield, respectively, the shapes
$(2,4),(1,1,4),(3,3)$. This example exhibits both varying width and a
nonsingleton preimage.

Every nonnegative capacity vector actually arises from a singleton
target. Set $s_t=2a_t+1$ and, backwards for $i<t$, set
\[
 s_i=2\left(a_i+\left\lceil \frac{s_{i+1}}2\right\rceil\right).
\]
The sequence $s$ is positive and weakly decreasing. Its entries other
than $s_t$ are even, so reversing it gives a $2$-singleton shape.
Substitution in \eqref{eq:capacity} recovers $a$.

\section{Discrete polynomial propagation}

For $t\geq1$ define the redundant individual bounds
\[
 U_1=a_1,\qquad U_i=\min(a_{i-1},a_i)\quad(2\leq i\leq t).
\]
Let $f_i(x)$ count prefixes $(d_1,\ldots,d_i)$ with last coordinate
$x$, all edges within the prefix satisfied, and
$0\leq d_j\leq U_j$ for every $j\leq i$. At an intermediate prefix
this anticipates the next edge, but every full solution already obeys
these bounds. Fixing the last coordinate gives
\begin{align}
 f_1(x)&=1 &&(0\leq x\leq U_1),\label{eq:initial}\\
 H_i(z)&=\sum_{x=0}^z f_i(x) &&(0\leq z\leq U_i),\label{eq:prefix}\\
 f_{i+1}(y)&=H_i(\min\{U_i,a_i-y\})
 &&(0\leq y\leq U_{i+1}).\label{eq:transition}
\end{align}
The argument in \eqref{eq:transition} is nonnegative since
$U_{i+1}\leq a_i$. The desired answer is $H_t(U_t)$.
For $t=1$ this is simply $a_1+1$.

Throughout, an interval $[l,r]$ means its integer points. On a
nonempty interval store a polynomial in the shifted binomial basis,
\begin{equation}\label{eq:piece}
 f_i(x)=\sum_{k=0}^{d}c_k\binom{x-l}{k},\qquad
 d\leq\min(i-1,r-l).
\end{equation}
Only values on the interval matter. In particular, a short interval
may use a lower-degree interpolant even if a polynomial previously
defined on a larger interval had higher degree.

\begin{lemma}\label{lem:pieces}
The recurrence \eqref{eq:initial}--\eqref{eq:transition} admits a
representation \eqref{eq:piece} on at most $i$ disjoint consecutive
nonempty intervals covering $[0,U_i]$. All coefficients are integers,
and they can be computed without iterating over a capacity-sized domain.
\end{lemma}
\begin{proof}
The first layer is $[0,U_1]$ with coefficient vector $(1)$.
Suppose a layer-$i$ piece is \eqref{eq:piece}, and let $T$ be the
total count on preceding intervals. The identity
\[
 \sum_{q=0}^n\binom qk=\binom{n+1}{k+1}
\]
and Pascal's identity show that the inclusive prefix on this interval
has coefficients, relative to $l$,
\begin{equation}\label{eq:prefixcoeff}
 h_0=T+c_0,\qquad h_k=c_{k-1}+c_k\ (1\leq k\leq d),
 \qquad h_{d+1}=c_d.
\end{equation}
Evaluating it at $r$ supplies the next value of $T$.
Thus prefix summation preserves the intervals and raises the degree
by at most one.

For each prefix interval $[l,r]$, reflection in
\eqref{eq:transition} gives the interval
\begin{equation}\label{eq:reflect}
 [a_i-r,a_i-l]\cap[0,U_{i+1}].
\end{equation}
Discard empty intersections. Each surviving interval carries the
reflected prefix polynomial, of degree at most $i$.
The only possible additional piece is the constant $H_i(U_i)$ on
\begin{equation}\label{eq:saturation}
 [0,\min\{U_{i+1},a_i-U_i-1\}],
\end{equation}
if this interval is nonempty. The strict endpoint is needed:
$a_i-y=U_i$ belongs to \eqref{eq:reflect}, whereas
\eqref{eq:saturation} accounts for $a_i-y>U_i$.
The resulting intervals are disjoint and cover $[0,U_{i+1}]$.
Reversing the old interval order puts the reflected pieces in
increasing order, so sorting is unnecessary.

On a surviving interval $[b,e]$ with reflected polynomial $Q$, set
\[
 D=\min\{i,\degp Q,e-b\}.
\]
Evaluate $Q$ at $b,b+1,\ldots,b+D$ and store their successive
forward differences. If $D=e-b$, every point of the interval was
sampled. Otherwise $\degp Q\leq D$, and Newton interpolation
recovers $Q$ identically. Hence the stored representation is exact,
integral, and has degree at most $\min(i,e-b)$.
Each old piece produces at most one new piece, and saturation adds
at most one. Induction proves the claim, including zero capacities,
one-point intervals and coincident boundaries.
\end{proof}

Equations \eqref{eq:prefixcoeff}--\eqref{eq:saturation}, followed by
the stated sampling rule, specify the algorithm completely: initialize
\eqref{eq:initial}, form the prefix pieces, reflect and normalize for
each next coordinate, and return the last prefix total. No symbolic
expansion in a monomial basis is required.

\section{Arithmetic, storage and bit bounds}

The normalization on short intervals also controls integer sizes.
The coefficient of \eqref{eq:piece} is
\begin{equation}\label{eq:diffbound}
 c_k=\Delta^k f_i(l)
 =\sum_{j=0}^k(-1)^{k-j}\binom kj f_i(l+j),\qquad
 |c_k|\leq 2^k(A+1)^{i-1}.
\end{equation}
Indeed all sampled arguments lie in the actual interval. Fixing the
last coordinate leaves at most $(A+1)^{i-1}$ prefixes, which proves
the bound. Moreover $0\leq T\leq(A+1)^i$, so the prefix coefficients
in \eqref{eq:prefixcoeff} have $O(i\log(A+2))$ bits. Intermediate
rows of the forward-difference table satisfy the same type of bound.

An evaluation at an offset $0\leq n\leq A$ uses
\[
 \binom n0=1,\qquad
 \binom n{k+1}=\binom nk\frac{n-k}{k+1}.
\]
Each division is exact. If $k>n$, the binomial value has already become
zero. For $k\leq t$, each nonzero binomial factor is at most
$(A+t+2)^t$. Combining this estimate with \eqref{eq:diffbound}
and \eqref{eq:prefixcoeff} bounds coefficient products, signed partial
sums, and temporary products before division by
$O(t\log(A+t+2))$ bits. This controls intermediate cancellation as
well as the final values. Endpoint arithmetic and counters fit in
the same bound.

At transition $i$ there are $O(i)$ pieces and $O(i)$ coefficients
per piece. Each of the $O(i)$ sample evaluations costs $O(i)$ integer
operations, and its forward-difference table costs $O(i^2)$ operations.
The transition therefore costs $O(i^3)$ operations. Prefix construction
and final summation cost less. Summing over the transitions yields
$O(t^4)$ operations. Keeping the current, prefix and next layers, and
one temporary difference table, takes $O(t^2)$ stored integers.
Schoolbook multiplication and division of $b$-bit integers cost
$O(b^2)$, giving the bit bound in Theorem~\ref{thm:main}.

Finally $\card{\mathcal D(a)}\leq(A+1)^t$, so the output has
$O(t\log(A+2))$ bits. For the target input,
$A\leq c_t/2$, $t=O(L)$ and $\log(A+2)=O(L)$.
Proposition~\ref{prop:bijection} and Lemma~\ref{lem:pieces} therefore
complete the proof of Theorem~\ref{thm:main}.

\section{Exact verification and limitations}

The accompanying Python implementation uses arbitrary-precision
integers throughout. A separate verifier constructs full inverse sets
by column-multiplicity enumeration, counts directly from
\eqref{eq:polytope}, and evaluates a suffix recurrence without
polynomial compression. Its exact finite coverage is recorded in
Table~\ref{tab:checks}. All these comparisons agreed.

\begin{table}[htbp]
\centering\small
\begin{tabular}{@{}p{0.32\linewidth}p{0.61\linewidth}@{}}
\toprule
Check & Coverage \\
\midrule
Capacity enumeration & All $5\,461$ vectors in
$\{0,1,2,3\}^t$, $0\leq t\leq6$ \\
Ferrers shapes and inverse sets & All $7\,338$ shapes of area at most
$24$, grouped into $1\,866$ fibres \\
Every interval value & $400$ seeded vectors, lengths $1$--$25$ and
capacities $0$--$40$; $74\,587$ count values \\
Large heterogeneous dilations & $16$ vectors, lengths $2$--$14$ and
$18$, at scales $10^{75}+31,\ldots,10^{75}+46$ \\
\bottomrule
\end{tabular}
\caption{Exact implementation checks. The unlimited theorem is proved
above; these finite comparisons test its implementation.}
\label{tab:checks}
\end{table}

For the large-dilation check, let $P(b)\subseteq\R^t$ be the real
polytope given by \eqref{eq:polytope} with capacities $b$.
Signing alternate columns of the edge-constraint matrix turns it
into the transpose of a directed incidence matrix. Such a matrix
is totally unimodular: in a square submatrix, expand along a row
with at most one nonzero entry; if every row has two, their opposite
signs make the matrix singular. Adjoining the terminal and
nonnegativity unit rows preserves this property by the same determinant
expansion. Hence every vertex of $P(b)$ is integral, by Cramer's rule.
It is bounded, and $P(kb)=kP(b)$ for $k\geq0$.
Ehrhart's theorem \cite[Theorem~3.8]{BR2015} makes
$\card{\mathcal D(kb)}$ a polynomial in $k$ of degree at most $t$,
also when zero capacities lower the dimension. Values at $0,\ldots,t$
from the suffix counter determine it by exact rational Lagrange
interpolation. One further small scale checks each interpolant before
evaluation at the large scale. This oracle shares neither the
binomial evaluation nor the finite-difference routine with the counter.

The files \texttt{count.py} and \texttt{verify.py} contain the counter
and its tests. The file \texttt{verification.json} records the inputs,
seeds and results. Running \texttt{python3 -B verify.py} from their
directory reproduces Table~\ref{tab:checks}. The functions
\texttt{fibre\_count(C)} and \texttt{compressed\_count(a)} return
exact integers. For example,
\texttt{python3 -B count.py '[1,1]'} reports the binary count
\texttt{0b11}.

Finite tests establish neither the unrestricted theorem nor priority.
The proof is specific to pairing rows: it provides no evaluation
theorem for arbitrary $m\geq3$, and it does not enumerate entire
$m$-level rook-equivalence classes.

\paragraph{Computational and AI assistance.}
OpenAI Codex language-model assistance was used in literature searches,
the development and checking of the argument, exact test programs,
and manuscript preparation. A separate model context reviewed the
argument and constructed additional verification routines; this was
neither external peer review nor proof-assistant certification.
The mathematical justification is the displayed proof. No empirical
data or floating-point computations are used.

\bibliographystyle{plainnat}
\bibliography{references}
\end{document}
